\newif\ifusectex
\IfFileExists{ctexart.cls}{%
  \usectextrue
  \documentclass[UTF8,12pt,a4paper]{ctexart}
}{%
  \usectexfalse
  \documentclass[12pt,a4paper]{article}
}

\ifusectex\else
  \usepackage{fontspec}
  \setmainfont[
    AutoFakeBold=2,
    AutoFakeSlant=0.2
  ]{Noto Sans CJK SC}
  \setsansfont[AutoFakeBold=2,AutoFakeSlant=0.2]{Noto Sans CJK SC}
  \setmonofont{DejaVu Sans Mono}
  \XeTeXlinebreaklocale "zh"
  \XeTeXlinebreakskip=0pt plus 1pt
  \renewcommand{\contentsname}{目录}
  \renewcommand{\abstractname}{摘要}
\fi

\usepackage[a4paper,margin=2.45cm,headheight=15pt]{geometry}
\usepackage{amsmath,amssymb,mathtools,bm}
\usepackage{booktabs,array,longtable}
\usepackage{xcolor}
\usepackage{enumitem}
\usepackage{fancyhdr}
\usepackage{hyperref}
\usepackage{listings}
\renewcommand{\refname}{参考文献}

\definecolor{sourcegray}{RGB}{78,88,102}
\definecolor{sourcebg}{RGB}{244,247,250}
\definecolor{linkblue}{RGB}{31,78,121}

\hypersetup{
  colorlinks=true,
  linkcolor=linkblue,
  citecolor=linkblue,
  urlcolor=linkblue,
  pdftitle={线性对流方程的有限体积半离散与全显式求解器推导},
  pdfauthor={张云飞}
}

\pagestyle{fancy}
\fancyhf{}
\lhead{线性对流方程的有限体积离散}
\rhead{参考书推导版}
\cfoot{\thepage}

\setlength{\parindent}{2em}
\setlength{\parskip}{0.25em}
\setlist[itemize]{leftmargin=2.2em,itemsep=0.25em,topsep=0.35em}
\setlist[enumerate]{leftmargin=2.4em,itemsep=0.35em,topsep=0.35em}
\numberwithin{equation}{section}

\newcommand{\Rone}{\textnormal{[R1]}}
\newcommand{\Rtwo}{\textnormal{[R2]}}
\newcommand{\Rthree}{\textnormal{[R3]}}
\newcommand{\Task}{\textnormal{[题目]}}
\newcommand{\source}[1]{%
  \par\smallskip
  \noindent\colorbox{sourcebg}{%
    \parbox{0.955\linewidth}{\footnotesize\color{sourcegray}\textbf{参考来源：}#1}%
  }
  \par\smallskip
}
\newcommand{\dd}{\,\mathrm d}
\newcommand{\p}{\partial}
\newcommand{\Sf}{\bm S}
\newcommand{\uf}{\bm u}
\newcommand{\phiv}{\bm\phi}

\lstdefinestyle{solver}{
  basicstyle=\ttfamily\small,
  frame=single,
  framerule=0.4pt,
  rulecolor=\color{sourcegray},
  backgroundcolor=\color{sourcebg},
  numbers=left,
  numberstyle=\tiny\color{sourcegray},
  stepnumber=1,
  numbersep=8pt,
  columns=fullflexible,
  keepspaces=true,
  showstringspaces=false,
  breaklines=true,
  xleftmargin=1.4em,
  framexleftmargin=1.0em
}

\title{\bfseries 线性对流方程的有限体积半离散\\与全显式数值求解器推导}
\author{张云飞}
\date{2026年8月22日}

\begin{document}

\maketitle

\begin{abstract}
本文针对守恒型线性对流方程
\[
  \frac{\partial\phi}{\partial t}+\nabla\cdot(\bm u\phi)=0
\]
给出严格依据指定参考书目的有限体积离散推导。全文分为四部分：第一部分从控制体积分、高斯定理和面高斯积分出发，得到时间连续的有限体积半离散形式；第二部分采用参考书介绍的一阶迎风面值，构造显式空间离散算子；第三部分从Taylor展开推导前向Euler时间格式，并根据参考书中的Courant数定义确定时间步；第四部分将网格数据、owner-neighbour通量装配、周期边界、时间循环和OpenFOAM显式算子对应关系组合成完整求解器。本文采用的一阶迎风和前向Euler是参考书明确给出的、能够构成完整全显式求解器的一组具体选择。
\end{abstract}

\tableofcontents
\newpage

\section*{符号约定与参考书对应关系}
\addcontentsline{toc}{section}{符号约定与参考书对应关系}

题目给出的控制方程为
\begin{equation}
  \frac{\partial\phi}{\partial t}
  +\nabla\cdot(\bm u\phi)=0.
  \label{eq:governing}
\end{equation}
其中，$\phi(\bm x,t)$ 为被输运的标量，$\bm u(\bm x,t)$ 为给定速度场。本题是参考书一般标量输运方程在密度取 $\rho=1$、扩散系数取零、源项取零时的特例。

\source{\Rone，第5.2节式(5.1)：一般标量输运方程；在该式中令 $\rho=1$、$\Gamma^\phi=0$、$Q^\phi=0$，即得到式\eqref{eq:governing}。\Task，第1.1节给出本题方程。}

为避免OpenFOAM中常用的面体积通量变量名 \texttt{phi} 与本题未知量 $\phi$ 混淆，本文始终用 $F_f$ 表示速度通过面 $f$ 的体积通量，用 $H_f$ 表示标量对流通量。

\begin{longtable}{>{\centering\arraybackslash}p{0.18\linewidth}p{0.70\linewidth}}
\toprule
符号 & 含义\\
\midrule
$\Omega_c$ & 第 $c$ 个控制体\\
$V_c=|\Omega_c|$ & 控制体面积或体积\\
$\mathcal F_c$ & 属于控制体 $c$ 的所有面组成的集合\\
$\bm n_{cf}$ & 面 $f$ 相对于控制体 $c$ 的单位外法向量\\
$\bm S_{cf}=|S_f|\bm n_{cf}$ & 面 $f$ 相对于控制体 $c$ 的有向面积向量\\
$O_f,N_f$ & 内部面 $f$ 的owner单元与neighbour单元\\
$F_f=\bm u_f\cdot\bm S_f$ & 面速度通量\\
$H_f=F_f\phi_f$ & 标量通过面的数值对流通量\\
$\mathcal Q_c$ & 控制体 $c$ 内所有面标量通量的未归一化代数和\\
$R_c=\mathcal Q_c/V_c$ & 控制体 $c$ 的体积归一化对流残差（离散散度）\\
\bottomrule
\end{longtable}

\source{\Rtwo，第1.3.1节及式(1.22)--(1.26)：面面积向量、面通量、owner-neighbour寻址与通量装配。}

\newpage
\section{第一部分：有限体积半离散形式}

\subsection{控制体划分与单元平均值}

将计算区域 $\Omega$ 划分为互不重叠的控制体：
\begin{equation}
  \Omega=\bigcup_{c=1}^{N_c}\Omega_c,
  \qquad
  \Omega_c^{\circ}\cap\Omega_d^{\circ}=\varnothing
  \quad(c\neq d).
  \label{eq:mesh_partition}
\end{equation}
对固定网格，$\Omega_c$ 及其体积 $V_c$ 均不随时间变化。定义控制体平均值
\begin{equation}
  \phi_c(t)
  =\frac{1}{V_c}\int_{\Omega_c}\phi(\bm x,t)\dd V.
  \label{eq:cell_average}
\end{equation}
于是
\begin{equation}
  \int_{\Omega_c}\phi(\bm x,t)\dd V=V_c\phi_c(t).
  \label{eq:cell_integral}
\end{equation}

\source{\Rone，第5.5节式(5.35)--(5.36)：固定控制体上的体积分以及控制体平均值定义；\Rtwo，第1.2节：单元中心有限体积变量与控制体离散。}

\subsection{在控制体上积分控制方程}

对式\eqref{eq:governing}在控制体 $\Omega_c$ 上积分：
\begin{equation}
  \int_{\Omega_c}\frac{\partial\phi}{\partial t}\dd V
  +\int_{\Omega_c}\nabla\cdot(\bm u\phi)\dd V=0.
  \label{eq:cell_integrated_pde}
\end{equation}
因为控制体固定，可以交换时间微分与空间积分：
\begin{align}
  \int_{\Omega_c}\frac{\partial\phi}{\partial t}\dd V
  &=\frac{\mathrm d}{\mathrm dt}\int_{\Omega_c}\phi\dd V
  \label{eq:time_exchange}\\
  &=\frac{\mathrm d}{\mathrm dt}(V_c\phi_c)
  =V_c\frac{\mathrm d\phi_c}{\mathrm dt}.
  \label{eq:time_cell_average}
\end{align}

\source{\Rone，第5.5节式(5.34)--(5.36)：非稳态项先作控制体积分；对于固定网格，$V_c$ 与控制体表面不随时间变化。}

\subsection{利用高斯散度定理转化对流项}

对空间散度项应用高斯散度定理：
\begin{equation}
  \int_{\Omega_c}\nabla\cdot(\bm u\phi)\dd V
  =\oint_{\partial\Omega_c}(\bm u\phi)\cdot\dd\bm S.
  \label{eq:gauss}
\end{equation}
若控制体边界由面集合 $\mathcal F_c$ 构成，则
\begin{equation}
  \oint_{\partial\Omega_c}(\bm u\phi)\cdot\dd\bm S
  =\sum_{f\in\mathcal F_c}
  \int_{S_f}(\bm u\phi)\cdot\bm n_{cf}\dd S.
  \label{eq:face_sum_exact}
\end{equation}
将式\eqref{eq:time_cell_average}与式\eqref{eq:face_sum_exact}代回式\eqref{eq:cell_integrated_pde}，得到精确控制体守恒关系
\begin{equation}
  V_c\frac{\mathrm d\phi_c}{\mathrm dt}
  +\sum_{f\in\mathcal F_c}
  \int_{S_f}(\bm u\phi)\cdot\bm n_{cf}\dd S=0.
  \label{eq:exact_balance}
\end{equation}

\source{\Rone，第5.2节式(5.3)--(5.4)：利用散度定理将体积分转化为表面积分；第5.2.1节式(5.8)：将控制体表面积分写成各面通量积分之和。}

\subsection{采用一点面中心高斯积分}

参考书对面通量采用高斯积分。对面 $f$，一般形式为
\begin{equation}
  \int_{S_f}\bm J\cdot\dd\bm S
  \approx
  \sum_{q=1}^{N_q}\omega_q
  (\bm J\cdot\bm n)_{f,q}|S_f|.
  \label{eq:face_gauss_general}
\end{equation}
采用一个位于面中心 $\bm x_f$ 的积分点，且权重 $\omega_1=1$，对于本题 $\bm J=\bm u\phi$，得到
\begin{align}
  \int_{S_f}(\bm u\phi)\cdot\bm n_{cf}\dd S
  &\approx(\bm u_f\phi_f)\cdot\bm n_{cf}|S_f|
  \label{eq:midpoint1}\\
  &=(\bm u_f\cdot\bm S_{cf})\phi_f.
  \label{eq:midpoint2}
\end{align}
其中
\begin{equation}
  \bm S_{cf}=|S_f|\bm n_{cf}.
  \label{eq:area_vector}
\end{equation}
定义面速度通量
\begin{equation}
  \boxed{F_{cf}=\bm u_f\cdot\bm S_{cf}},
  \label{eq:face_flux_local}
\end{equation}
则
\begin{equation}
  \int_{S_f}(\bm u\phi)\cdot\bm n_{cf}\dd S
  \approx F_{cf}\phi_f.
  \label{eq:flux_approx}
\end{equation}

\source{\Rone，第5.2.1节式(5.11)--(5.12)：面高斯积分及单积分点面中心近似；\Rtwo，式(1.16)--(1.17)：面中心平均与散度项的面求和；\Rtwo，式(1.22)：$F_f=\bm U_f\cdot\bm S_f$。}

\subsection{半离散方程}

将式\eqref{eq:flux_approx}代入精确平衡式\eqref{eq:exact_balance}：
\begin{equation}
  V_c\frac{\mathrm d\phi_c}{\mathrm dt}
  +\sum_{f\in\mathcal F_c}F_{cf}\phi_f=0.
  \label{eq:semidiscrete_integral}
\end{equation}
除以控制体体积 $V_c$：
\begin{equation}
  \boxed{
  \frac{\mathrm d\phi_c}{\mathrm dt}
  =-\frac{1}{V_c}\sum_{f\in\mathcal F_c}F_{cf}\phi_f
  }.
  \label{eq:semidiscrete_general}
\end{equation}
定义体积归一化的离散空间残差
\begin{equation}
  R_c(\phiv)
  :=\frac{1}{V_c}\sum_{f\in\mathcal F_c}F_{cf}\phi_f,
  \label{eq:spatial_operator_general}
\end{equation}
则
\begin{equation}
  \boxed{
  \frac{\mathrm d\phi_c}{\mathrm dt}+R_c(\phiv)=0
  }.
  \label{eq:mol_cell}
\end{equation}
对全部控制体组装：
\begin{equation}
  \boxed{
  \frac{\mathrm d\phiv}{\mathrm dt}+\bm R_h(\phiv)=\bm0
  }.
  \label{eq:mol_global}
\end{equation}
式\eqref{eq:semidiscrete_general}--\eqref{eq:mol_global}中，空间散度已经转化为有限个面通量之和并除以控制体体积，而时间导数仍保持连续，因此它们就是题目要求的有限体积半离散形式。此时面值 $\phi_f$ 尚待空间插值格式封闭。

\source{\Rone，第5.2.3节式(5.15)：一点积分下的有限体积半离散通量和；第5.5节式(5.34)--(5.38)：保留时间项的非稳态有限体积关系；\Rtwo，式(1.23)：离散散度算子是全部控制体面通量之和。}

\newpage
\section{第二部分：显式空间离散}

\subsection{“显式空间离散”的参考书含义}

参考书将OpenFOAM有限体积算子分为两类：\texttt{fvc::} 直接依据当前场值返回离散微分算子的场，是显式离散；\texttt{fvm::} 返回由对角、上三角、下三角系数及源项构成的矩阵，是隐式离散。本题要求显式空间离散，因此空间残差必须由当前已知的单元值直接计算，不将新的未知时间层组装进线性方程组。

\source{\Rone，第5.10.2节表5.1：\texttt{fvc::} 对应显式微分算子，\texttt{fvm::} 对应隐式微分算子；第5.10.2节随后对 \texttt{fvc::div} 与 \texttt{fvm::div} 返回对象的区别作出说明。}

为了得到一个完整、明确且与参考书稳定性推导一致的空间算子，本文选择参考书的一阶迎风格式计算 $\phi_f$。

\source{\Rone，第11.2.4节：一阶迎风格式；\Rthree，第59.2.1节：迎风格式以面上游单元的单元中心值作为面值。}

\subsection{owner-neighbour面方向与面通量}

对任意内部面 $f$，OpenFOAM只定义一个面积向量 $\bm S_f$，其方向从owner单元 $O_f$ 指向neighbour单元 $N_f$。定义
\begin{equation}
  F_f=\bm u_f\cdot\bm S_f.
  \label{eq:global_face_flux}
\end{equation}
$F_f$ 的符号给出相对于该唯一面方向的输运方向：
\begin{itemize}
  \item $F_f>0$：流动从 $O_f$ 指向 $N_f$；
  \item $F_f<0$：流动从 $N_f$ 指向 $O_f$。
\end{itemize}

\source{\Rtwo，式(1.22)：标量体积面通量；式(1.24)--(1.25)：owner与neighbour索引；式(1.26)前后的文字：$\bm S_f$ 从owner指向neighbour。}

\subsection{一阶迎风面值}

迎风原则是：面值取信息来源方向上的上游单元值。因此
\begin{equation}
  \boxed{
  \phi_f^{\mathrm{up}}=
  \begin{cases}
    \phi_{O_f}, & F_f\ge0,\\[1mm]
    \phi_{N_f}, & F_f<0.
  \end{cases}
  }
  \label{eq:upwind_piecewise}
\end{equation}
定义正、负通量分裂
\begin{equation}
  F_f^+:=\max(F_f,0),
  \qquad
  F_f^-:=\min(F_f,0).
  \label{eq:flux_split}
\end{equation}
则数值对流通量可以统一写成
\begin{equation}
  \boxed{
  H_f:=F_f\phi_f^{\mathrm{up}}
  =F_f^+\phi_{O_f}+F_f^-\phi_{N_f}
  }.
  \label{eq:upwind_flux_global}
\end{equation}
确实，当 $F_f\ge0$ 时有 $F_f^+=F_f,F_f^-=0$，式\eqref{eq:upwind_flux_global}退化为 $H_f=F_f\phi_{O_f}$；当 $F_f<0$ 时有 $F_f^+=0,F_f^-=F_f$，它退化为 $H_f=F_f\phi_{N_f}$。

\source{\Rone，第11.2.4节：迎风面值依赖流动方向并取上游节点值；\Rthree，第59.2.1节：迎风格式的面值定义。式\eqref{eq:flux_split}--\eqref{eq:upwind_flux_global}是对参考书分段迎风定义的等价代数写法。}

\subsection{内部面的守恒装配}

对内部面 $f$ 只计算一次 $H_f$。在程序装配阶段，先把它加入未除以体积的面通量累加量 $\mathcal Q_c$：
\begin{align}
  \mathcal Q_{O_f}&\mathrel{+}=H_f,
  \label{eq:owner_add}\\
  \mathcal Q_{N_f}&\mathrel{-}=H_f.
  \label{eq:neighbour_subtract}
\end{align}
同一面通量在两个相邻控制体方程中的大小相等、符号相反，这正是局部守恒的离散实现。全部面装配完成后，再对每个控制体归一化：
\begin{equation}
  R_c(\phiv)=\frac{\mathcal Q_c(\phiv)}{V_c}.
  \label{eq:residual_from_accumulator}
\end{equation}
因此，$\mathcal Q_c$ 是装配过程中的中间量，$R_c$ 才是本文统一使用的最终残差。

\source{\Rtwo，式(1.26)：每个内部面的离散散度贡献对owner取加号、对neighbour取减号；\Rone，第5.2.1节式(5.8)--(5.10)后的说明：面通量形式使有限体积法具有守恒性。}

\subsection{封闭的迎风半离散方程}

现在改用相对于每个控制体 $c$ 的外法向量。令 $d(f)$ 表示面 $f$ 另一侧的相邻控制体，并定义
\begin{equation}
  F_{cf}=\bm u_f\cdot\bm S_{cf}.
  \label{eq:local_flux_again}
\end{equation}
则迎风面值为
\begin{equation}
  \phi_{cf}^{\mathrm{up}}=
  \begin{cases}
    \phi_c, & F_{cf}\ge0,\\[1mm]
    \phi_{d(f)}, & F_{cf}<0,
  \end{cases}
  \label{eq:upwind_local}
\end{equation}
并且
\begin{equation}
  F_{cf}\phi_{cf}^{\mathrm{up}}
  =F_{cf}^+\phi_c+F_{cf}^-\phi_{d(f)}.
  \label{eq:local_split_flux}
\end{equation}
将式\eqref{eq:local_split_flux}代入一般半离散式\eqref{eq:semidiscrete_general}，得到封闭的一阶迎风空间半离散方程
\begin{equation}
  \boxed{
  \frac{\mathrm d\phi_c}{\mathrm dt}
  =-\frac{1}{V_c}
  \sum_{f\in\mathcal F_c}
  \left(F_{cf}^+\phi_c+F_{cf}^-\phi_{d(f)}\right)
  }.
  \label{eq:upwind_semidiscrete}
\end{equation}
定义显式迎风空间残差
\begin{equation}
  R_c(\phiv)
  :=\frac{1}{V_c}\sum_{f\in\mathcal F_c}
  \left(F_{cf}^+\phi_c+F_{cf}^-\phi_{d(f)}\right),
  \label{eq:residual}
\end{equation}
则
\begin{equation}
  \boxed{
  \frac{\mathrm d\phi_c}{\mathrm dt}
  =-R_c(\phiv)
  }.
  \label{eq:residual_ode}
\end{equation}

\source{\Rone，第5.2.3节：面通量代入有限体积方程；第11.2.4节：迎风面值；\Rtwo，式(1.23)与(1.26)：面插值后的离散散度及owner-neighbour装配。式\eqref{eq:upwind_semidiscrete}是这些参考公式针对本题纯对流方程的直接组合。}

\subsection{离散质量守恒}

对全部控制体的式\eqref{eq:residual_ode}乘以 $V_c$ 后求和：
\begin{equation}
  \frac{\mathrm d}{\mathrm dt}\sum_cV_c\phi_c
  =-\sum_cV_cR_c.
  \label{eq:global_mass1}
\end{equation}
由于 $V_cR_c=\mathcal Q_c$，内部面按照式\eqref{eq:owner_add}--\eqref{eq:neighbour_subtract}成对抵消，因此只剩边界通量：
\begin{equation}
  \frac{\mathrm d}{\mathrm dt}\sum_cV_c\phi_c
  =-\sum_{f\subset\partial\Omega}H_f.
  \label{eq:global_mass2}
\end{equation}
对于题目正弦波算例规定的周期边界，周期配对面的通量同样大小相等、符号相反，于是
\begin{equation}
  \boxed{
  \frac{\mathrm d}{\mathrm dt}\sum_cV_c\phi_c=0
  }.
  \label{eq:mass_conservation}
\end{equation}

\source{\Rone，第5.2.1节关于面通量守恒性的说明；\Rtwo，式(1.26)的owner-neighbour反号装配；\Task，第1.2.1节规定 $x$、$y$ 方向采用周期边界。}

\newpage
\section{第三部分：显式时间离散}

\subsection{前向Euler格式的Taylor推导}

空间离散完成后，方程成为常微分方程
\begin{equation}
  \frac{\mathrm d\phi_c}{\mathrm dt}
  =-R_c(\phiv).
  \label{eq:ode_before_time}
\end{equation}
在时刻 $t^n$ 对 $\phi_c(t^n+\Delta t)$ 作Taylor展开：
\begin{equation}
  \phi_c(t^n+\Delta t)
  =\phi_c(t^n)
  +\Delta t\left.\frac{\mathrm d\phi_c}{\mathrm dt}\right|_{t^n}
  +\frac{\Delta t^2}{2}
  \left.\frac{\mathrm d^2\phi_c}{\mathrm dt^2}\right|_{t^n}
  +O(\Delta t^3).
  \label{eq:taylor_time}
\end{equation}
记
\begin{equation}
  \phi_c^n=\phi_c(t^n),
  \qquad
  \phi_c^{n+1}=\phi_c(t^n+\Delta t),
  \label{eq:time_notation}
\end{equation}
由式\eqref{eq:taylor_time}移项并除以 $\Delta t$：
\begin{equation}
  \left.\frac{\mathrm d\phi_c}{\mathrm dt}\right|_{t^n}
  =\frac{\phi_c^{n+1}-\phi_c^n}{\Delta t}+O(\Delta t).
  \label{eq:forward_difference}
\end{equation}
忽略 $O(\Delta t)$ 项，得到一阶前向Euler近似
\begin{equation}
  \left.\frac{\mathrm d\phi_c}{\mathrm dt}\right|_{t^n}
  \approx\frac{\phi_c^{n+1}-\phi_c^n}{\Delta t}.
  \label{eq:forward_euler_derivative}
\end{equation}

\source{\Rone，第13.2.1节式(13.5)--(13.6)：在时间 $t$ 处向前Taylor展开，并得到一阶前向差分。}

\subsection{将空间算子取在旧时间层}

前向Euler要求空间算子全部在旧时间层 $t^n$ 计算。将式\eqref{eq:forward_euler_derivative}代入式\eqref{eq:ode_before_time}：
\begin{equation}
  \frac{\phi_c^{n+1}-\phi_c^n}{\Delta t}
  =-R_c(\phiv^n).
  \label{eq:euler_residual}
\end{equation}
解出新时间层：
\begin{equation}
  \boxed{
  \phi_c^{n+1}
  =\phi_c^n-\Delta t R_c(\phiv^n)
  }.
  \label{eq:euler_update_residual}
\end{equation}

\source{\Rone，第13.2.1节式(13.7)：瞬态项采用前向差分，空间算子 $L(\phi_C^t)$ 取在旧时间 $t$；该节指出此时无需求解方程组。}

\subsection{完全离散的一阶迎风-前向Euler格式}

将已经除以控制体体积的显式迎风残差式\eqref{eq:residual}代入式\eqref{eq:euler_update_residual}：
\begin{equation}
  \boxed{
  \phi_c^{n+1}
  =\phi_c^n
  -\frac{\Delta t}{V_c}
  \sum_{f\in\mathcal F_c}
  \left[
    (F_{cf}^n)^+\phi_c^n
    +(F_{cf}^n)^-\phi_{d(f)}^n
  \right]
  }.
  \label{eq:fully_discrete}
\end{equation}
式\eqref{eq:fully_discrete}右端只含旧时间层的 $\phi_c^n$、$\phi_{d(f)}^n$ 和 $F_{cf}^n$，所以可直接计算 $\phi_c^{n+1}$，无需组装或求解线性代数方程组。

\source{\Rone，第13.2.1节式(13.7)--(13.12)：显式Euler在旧时间层计算全部空间项；\Rone，第11.2.4节：迎风面值。式\eqref{eq:fully_discrete}是两部分参考公式的直接代入结果。}

\subsection{Courant数与时间步}

参考手册将单元面通量绝对值之和记为
\begin{equation}
  \mathrm{sumPhi}_c
  =\sum_{f\in\mathcal F_c}|F_{cf}|.
  \label{eq:sumphi}
\end{equation}
单元Courant数为
\begin{equation}
  Co_c
  =\frac{\Delta t}{2V_c}
  \sum_{f\in\mathcal F_c}|F_{cf}|,
  \label{eq:local_co}
\end{equation}
全局最大Courant数为
\begin{equation}
  Co_{\max}
  =\frac{\Delta t}{2}
  \max_c\left(
    \frac{\sum_{f\in\mathcal F_c}|F_{cf}|}{V_c}
  \right).
  \label{eq:max_co}
\end{equation}
题目规定所有计算采用
\begin{equation}
  Co_{\max}=0.2.
  \label{eq:co_task}
\end{equation}
令式\eqref{eq:max_co}等于 $0.2$，得到时间步
\begin{equation}
  \boxed{
  \Delta t
  =\frac{2\times0.2}{
    \displaystyle
    \max_c\left(
      \frac{\sum_{f\in\mathcal F_c}|F_{cf}|}{V_c}
    \right)}
  =\frac{0.4}{
    \displaystyle
    \max_c\left(
      \frac{\sum_{f\in\mathcal F_c}|F_{cf}|}{V_c}
    \right)}
  }.
  \label{eq:dt_formula}
\end{equation}
若速度场随时间变化，则 $F_{cf}=F_{cf}(t)$，需要在每个时间步重新计算式\eqref{eq:dt_formula}。

\source{\Rthree，第57.6.4节式(171)--(174)：\texttt{sumPhi}、最大Courant数及面通量绝对值之和的有限体积表达；\Task，第1.2.1节规定 $CFL=0.2$。}

\subsection{迎风-前向Euler格式的非负系数条件}

将式\eqref{eq:fully_discrete}整理为
\begin{align}
  \phi_c^{n+1}
  =&\left[
    1-\frac{\Delta t}{V_c}
    \sum_f(F_{cf}^n)^+
  \right]\phi_c^n
  \nonumber\\
  &+\sum_f\left[
    -\frac{\Delta t}{V_c}(F_{cf}^n)^-
  \right]\phi_{d(f)}^n.
  \label{eq:convex_form}
\end{align}
由于 $(F_{cf}^n)^-\le0$，邻居系数非负。为了使中心系数也非负，需要
\begin{equation}
  \frac{\Delta t}{V_c}\sum_f(F_{cf}^n)^+\le1.
  \label{eq:positive_condition}
\end{equation}
对于本题两个算例的无散速度场，离散面通量满足
\begin{equation}
  \sum_fF_{cf}=0.
  \label{eq:discrete_div_free}
\end{equation}
于是
\begin{equation}
  \sum_fF_{cf}^+
  =\frac12\sum_f|F_{cf}|,
  \label{eq:positive_half_abs}
\end{equation}
式\eqref{eq:positive_condition}化为
\begin{equation}
  \frac{\Delta t}{2V_c}\sum_f|F_{cf}|\le1,
  \qquad\text{即}\qquad Co_c\le1.
  \label{eq:co_stability}
\end{equation}
题目采用 $Co=0.2$，位于该显式迎风稳定限制之内。

\source{\Rone，第13.2.2节式(13.13)--(13.18)：以前向Euler和迎风离散一维瞬态对流问题，并由相反符号法则得到 $CFL_{\mathrm{conv}}\le1$。式\eqref{eq:convex_form}--\eqref{eq:co_stability}是同一条件在一般控制体面通量记号下的等价推导。}

\newpage
\section{第四部分：完整全显式数值求解器}

\subsection{求解器输入与离散数据}

求解器需要以下输入：
\begin{enumerate}
  \item 控制体体积 $V_c$ 与中心 $\bm x_c$；
  \item 每个面的中心 $\bm x_f$ 与面积向量 $\bm S_f$；
  \item 每个内部面的owner索引 $O_f$ 与neighbour索引 $N_f$；
  \item 边界面的类型及周期匹配关系；
  \item 给定速度函数 $\bm u(\bm x,t)$；
  \item 初始条件 $\phi_0(\bm x)$、终止时间 $T$ 和目标Courant数 $0.2$。
\end{enumerate}

这些数据对三角形和四边形网格使用同一套计算公式；网格类型只改变每个控制体拥有的面数、体积和面面积向量，不改变通量装配结构。

\source{\Rtwo，第1.3节及式(1.24)--(1.26)：一般多面体有限体积网格的面、owner和neighbour寻址；\Task，第1.2节规定三角形、四边形网格、速度场、终止时间和CFL。}

\subsection{初始条件离散}

严格的有限体积初值是初始函数的控制体平均：
\begin{equation}
  \boxed{
  \phi_c^0
  =\frac{1}{V_c}\int_{\Omega_c}\phi_0(\bm x)\dd V
  }.
  \label{eq:initial_average}
\end{equation}
若采用与基础有限体积离散一致的单元中心一点积分，则
\begin{equation}
  \phi_c^0\approx\phi_0(\bm x_c).
  \label{eq:initial_midpoint}
\end{equation}
若后续更换为高阶空间格式并测量高阶收敛率，则初值体积分也应采用足够高阶的体积分公式，避免初值积分误差限制总体精度。

\source{\Rone，第5.5节式(5.36)：控制体平均值；第5.2.2节式(5.14)：体积分的高斯积分；第5.2.3节说明基础有限体积法通常采用一个积分点。}

\subsection{单个时间步的严格执行顺序}

假设时刻 $t^n$ 的全部单元值 $\phiv^n$ 已知。一个时间步必须按如下顺序完成。

\subsubsection*{步骤1：冻结旧时间层并清零残差}

保存
\begin{equation}
  \phi_c^{\mathrm{old}}=\phi_c^n,
  \label{eq:freeze_old}
\end{equation}
并令
\begin{equation}
  \mathcal Q_c^n=0,
  \qquad c=1,\ldots,N_c.
  \label{eq:zero_residual}
\end{equation}
这里清零的是未除以体积的面通量累加量；全部面装配完成后再定义
\begin{equation}
  R_c^n=\frac{\mathcal Q_c^n}{V_c}.
  \label{eq:solver_residual_normalization}
\end{equation}
在全部面通量计算结束之前，不修改 $\phi_c^n$。

\source{\Rone，第13.2.1节式(13.7)及其说明：显式Euler的全部空间项必须在旧时间 $t$ 计算。}

\subsubsection*{步骤2：计算内部面通量}

对每个内部面 $f$：
\begin{align}
  \bm u_f^n&=\bm u(\bm x_f,t^n),
  \label{eq:solver_uf}\\
  F_f^n&=\bm u_f^n\cdot\bm S_f,
  \label{eq:solver_F}\\
  \phi_f^n&=
  \begin{cases}
    \phi_{O_f}^n,&F_f^n\ge0,\\
    \phi_{N_f}^n,&F_f^n<0,
  \end{cases}
  \label{eq:solver_face_value}\\
  H_f^n&=F_f^n\phi_f^n.
  \label{eq:solver_H}
\end{align}
然后装配
\begin{align}
  \mathcal Q_{O_f}^n&\mathrel{+}=H_f^n,
  \label{eq:solver_owner}\\
  \mathcal Q_{N_f}^n&\mathrel{-}=H_f^n.
  \label{eq:solver_neighbour}
\end{align}

\source{\Rtwo，式(1.22)：$F_f=\bm U_f\cdot\bm S_f$；式(1.26)：owner加通量、neighbour减通量；\Rone，第11.2.4节与\Rthree，第59.2.1节：一阶迎风面值。}

\subsubsection*{步骤3：处理边界面}

对于周期边界，将一个周期面的上游状态从其周期匹配面另一侧的单元取得，并以相同的 $F_f\phi_f^{\mathrm{up}}$ 形式计算通量。周期成对通量需要保持大小相等、方向相反，从而维持式\eqref{eq:mass_conservation}的全局守恒。

\source{\Task，第1.2.1节明确规定 $x$、$y$ 方向周期边界；\Rone，第5.2.1节的面通量守恒原则；\Rtwo，式(1.26)的成对反号通量装配。}

\subsubsection*{步骤4：由Courant数计算时间步}

计算每个控制体的面通量绝对值之和，并通过式\eqref{eq:dt_formula}得到
\begin{equation}
  \Delta t^n
  =\frac{0.4}{
    \displaystyle
    \max_c\left(
      \frac{\sum_{f\in\mathcal F_c}|F_{cf}^n|}{V_c}
    \right)}.
  \label{eq:solver_dt}
\end{equation}
为使最后一步精确到达终止时间 $T$，采用
\begin{equation}
  \Delta t^n
  \leftarrow
  \min(\Delta t^n,T-t^n).
  \label{eq:last_step}
\end{equation}

\source{\Rthree，第57.6.4节式(171)--(174)：基于面通量与控制体体积的Courant数；\Task，第1.2.1节：$CFL=0.2$。式\eqref{eq:last_step}是使离散时间终点与题目指定输出时刻一致的时间循环实现。}

\subsubsection*{步骤5：统一执行前向Euler更新}

全部面通量和残差计算完毕后，对所有控制体同时更新：
\begin{equation}
  \boxed{
  \phi_c^{n+1}
  =\phi_c^n-\Delta t^n R_c^n
  }.
  \label{eq:solver_update}
\end{equation}
推进时间
\begin{equation}
  t^{n+1}=t^n+\Delta t^n.
  \label{eq:advance_time}
\end{equation}
必须先完成所有 $R_c^n$ 的计算，再统一更新 $\phi_c^{n+1}$；若在面循环中直接覆盖 $\phi_c^n$，后续面会混用新旧时间层，便不再是式(13.7)定义的前向Euler格式。

\source{\Rone，第13.2.1节式(13.7)：新时间层由旧时间层空间算子显式计算；式(13.12)：显式关系不求解方程组。}

\subsection{完整求解器伪代码}

\begin{lstlisting}[style=solver]
Input: mesh geometry, owner/neighbour addressing,
       boundary connectivity, velocity u(x,t), initial field phi0,
       final time T, target Co = 0.2

for every cell c:
    phi[c] = cellAverage(phi0 over cell c)

t = 0
while t < T:
    phiOld = phi
    fluxSum[c] = 0 for all cells

    for every internal face f:
        O = owner[f]
        N = neighbour[f]
        uFace = u(faceCentre[f], t)
        F = dot(uFace, faceAreaVector[f])

        if F >= 0:
            phiFace = phiOld[O]
        else:
            phiFace = phiOld[N]

        H = F * phiFace
        fluxSum[O] += H
        fluxSum[N] -= H

    assemble boundary and periodic-face fluxes
    residual[c] = fluxSum[c] / V[c] for all cells

    dt = 0.4 / max_c(sum_f(abs(F_cf)) / V[c])
    dt = min(dt, T - t)

    for every cell c:
        phi[c] = phiOld[c] - dt * residual[c]

    t = t + dt

Output: phi at t = T
\end{lstlisting}

\source{伪代码第13--27行对应\Rtwo 式(1.22)与(1.26)以及\Rthree 第59.2.1节；第30--31行对应\Rthree 第57.6.4节式(171)--(174)和\Task 的 $CFL=0.2$；第33--36行对应\Rone 第13.2.1节式(13.7)。}

\subsection{OpenFOAM显式实现对应关系}

在OpenFOAM的有限体积算子记号中，若表面标量场 $F$ 存储 $\bm u_f\cdot\bm S_f$，体标量场 $\phi$ 存储控制体值，则显式空间散度对应
\begin{equation}
  R_h(\phi)=\texttt{fvc::div(F, phi)}.
  \label{eq:fvcdiv}
\end{equation}
这里的 $R_h$ 已经包含控制体体积归一化，即
\begin{equation}
  R_{h,c}(\phi)
  =\frac{1}{V_c}\sum_{f\in\mathcal F_c}F_{cf}\phi_f.
  \label{eq:fvcdiv_normalized}
\end{equation}
若在离散格式字典中为该散度项选择 \texttt{Gauss upwind}，则 \texttt{fvc::div} 通过一阶迎风面值形成当前时间层的体积归一化离散散度场。前向Euler更新的概念表达为
\begin{equation}
  \boxed{
  \phi^{n+1}
  =\phi^n-\Delta t\,\texttt{fvc::div(F, phi}^{n}\texttt{)}
  }.
  \label{eq:openfoam_update}
\end{equation}
该实现不应将对流项写成 \texttt{fvm::div(F,phi)} 后调用矩阵求解，因为那会形成隐式空间离散，而不再是题目要求的全显式求解器。

\source{\Rone，第5.10.2节表5.1及其后说明：\texttt{fvc::} 返回显式计算场，\texttt{fvm::} 返回隐式系数矩阵；\Rone，第11.10节关于 \texttt{fvc::div} 显式面通量散度；\Rthree，第59.2.1节：\texttt{upwind} 面值。}

\subsection{求解器必须满足的基本检查}

\paragraph{常数保持。}
若 $\phi_c^n=C$ 且离散速度场满足 $\sum_fF_{cf}=0$，则
\begin{equation}
  R_c^n=\frac{C}{V_c}\sum_fF_{cf}=0,
\end{equation}
从而 $\phi_c^{n+1}=C$。这检查面方向与速度通量是否一致。

\paragraph{内部面守恒。}
对每个内部面，应满足owner贡献与neighbour贡献之和为零：
\begin{equation}
  H_f+(-H_f)=0.
\end{equation}
这检查owner-neighbour装配符号。

\paragraph{周期总量守恒。}
周期边界条件下，应有
\begin{equation}
  \sum_cV_c\phi_c^{n+1}
  =\sum_cV_c\phi_c^n.
\end{equation}
至舍入误差范围内成立。

\source{\Rone，第5.2.1节关于有限体积面通量守恒性的说明；\Rtwo，式(1.26)：内部面反号装配；\Task，第1.2.1节：周期边界。上述三项是对相应参考公式在程序中的直接一致性检查。}

\section*{最终公式汇总}
\addcontentsline{toc}{section}{最终公式汇总}

有限体积半离散形式为
\begin{equation}
  \boxed{
  \frac{\mathrm d\phi_c}{\mathrm dt}
  =-\frac{1}{V_c}\sum_{f\in\mathcal F_c}F_{cf}\phi_f,
  \qquad
  F_{cf}=\bm u_f\cdot\bm S_{cf}
  }.
\end{equation}

一阶迎风显式空间离散为
\begin{equation}
  \boxed{
  \frac{\mathrm d\phi_c}{\mathrm dt}
  =-\frac{1}{V_c}
  \sum_{f\in\mathcal F_c}
  \left(F_{cf}^+\phi_c+F_{cf}^-\phi_{d(f)}\right)
  }.
\end{equation}

一阶迎风-前向Euler完全离散格式为
\begin{equation}
  \boxed{
  \phi_c^{n+1}
  =\phi_c^n
  -\frac{\Delta t}{V_c}
  \sum_{f\in\mathcal F_c}
  \left[(F_{cf}^n)^+\phi_c^n
  +(F_{cf}^n)^-\phi_{d(f)}^n\right]
  }.
\end{equation}

题目指定 $Co=0.2$ 时的时间步为
\begin{equation}
  \boxed{
  \Delta t
  =\frac{0.4}{
    \displaystyle
    \max_c\left(
      \frac{\sum_{f\in\mathcal F_c}|F_{cf}|}{V_c}
    \right)}
  }.
\end{equation}

\source{半离散公式：\Rone 第5.2、5.2.1、5.5节与\Rtwo 式(1.17)、(1.22)--(1.23)；迎风公式：\Rone 第11.2.4节与\Rthree 第59.2.1节；前向Euler：\Rone 第13.2.1节式(13.5)--(13.7)；CFL时间步：\Rthree 第57.6.4节式(171)--(174)与\Task 的 $CFL=0.2$。}

\newpage
\begin{thebibliography}{9}

\bibitem{R1}
F. Moukalled, L. Mangani, and M. Darwish,
\textit{The Finite Volume Method in Computational Fluid Dynamics: An Advanced Introduction with OpenFOAM and Matlab},
Springer, 2016.

本文使用位置：第5.2节“半离散方程”、第5.2.1节“控制体面通量积分”、第5.5节“瞬态半离散方程”、第5.10.2节“OpenFOAM计算实现”、第11.2.4节“迎风格式”、第13.2.1节“前向Euler格式”和第13.2.2节“前向Euler稳定性”。

\bibitem{R2}
T. Mari\'{c}, J. H\"opken, and K. Mooney,
\textit{The OpenFOAM Technology Primer}, 2021 edition.

本文使用位置：式(1.16)--(1.17)面中心积分，式(1.22)面体积通量，式(1.23)离散散度，式(1.24)--(1.26) owner-neighbour寻址与通量装配。

\bibitem{R3}
G. Holzinger,
\textit{OpenFOAM: A Little User-Manual},
23 March 2020.

本文使用位置：第57.6.4节Courant数，式(168)--(174)；第59.2.1节upwind空间插值格式。

\end{thebibliography}

\end{document}
